\section{A Criterion for When Admission Control Can Pay}
\label{sec:criterion}

Throughput during busy time is
$\overline{R}/\tau$ \cite{denning1978}, so raising concurrency by a factor $f$
buys a factor $f^{1-\epsilon}$, where
\begin{equation}
  \epsilon \;=\; \frac{\mathrm{d}\ln \tau}{\mathrm{d}\ln \overline{R}}
  \label{eq:elasticity}
\end{equation}
is the elasticity of per-step cost with respect to the running count. At
$\epsilon = 0$ the step cost is flat, concurrency converts one for one into
throughput, and admission control has room; at $\epsilon = 1$ the two cancel
exactly and it has none.

On the five device-only configurations, where \S\ref{sec:model} establishes
that the step cost is driven by recomputation, $\overline{R}_{\text{busy}}$
runs from 1.00 to 1.84. A least-squares fit of Eq.~\ref{eq:elasticity} over
them gives $\epsilon = 0.61$ on the per-request measurement
($R^2 = 0.89$), or $0.38$ on the model-consistent one (Fig.~\ref{fig:cancel}).
In the units a controller works in, adding 10\%
concurrency buys 3.8\% to 6.1\% of busy-time throughput, against the 10\% it
would buy at $\epsilon = 0$.

The segment elasticities are $\epsilon = -0.21$ from \texttt{cap1} to
\texttt{cap2}, $+0.43$ from \texttt{cap2} to \texttt{cap3}, $+1.15$ from
\texttt{cap3} to \texttt{cap4}, and $+1.14$ from \texttt{cap4} to
\texttt{cap5}. Below $\overline{R}_{\text{busy}} \approx 1.5$ added
concurrency still buys something; above, $\epsilon$ passes one and busy-time
throughput stops rising with the running count. That is not the same as the run
getting slower end to end, since less idle time can more than pay for a longer
busy time. The single number $0.38$ averages those regimes and describes none
of them.

The controllers run in a different region: all nine controller arms run with
HiCache enabled, which the five device-only caps do not. (Every elasticity
here, fitted or segmented, uses $\overline{R}_{\text{busy}}$ at the two-decimal
precision of \texttt{decomp.csv}; at six decimals the \emph{fitted} values read
$0.37$ and $0.60$, and the segments move by $0.01$--$0.03$.) On the three
HiCache configurations that span the controllers' range, the same elasticity
gives $\epsilon = +0.19$ from \texttt{cap2}+HC to \texttt{cap3}+HC and $+0.46$
from \texttt{cap3}+HC to \texttt{cap4}+HC, over $\overline{R}_{\text{busy}}$
1.25 to 1.73---the range covering the whole controller operating region (1.35 to
1.69), over which $\epsilon$ stays below one. \texttt{cap5}+HC then adds a
segment with $\epsilon = +2.87$ whose denominator is near zero: its running
count moves only from 1.73 to 1.75 while the step cost moves by 1.0\,ms, and
\texttt{cap6}+HC at $\overline{R}_{\text{busy}}$ 1.87 gives $\epsilon = +0.48$.
We do \emph{not} read the $+2.87$ segment as ``the spike was an artifact'': both
step-cost measurements agree something real happens there, but the small
denominator costs us \emph{resolution}, not the existence of a cost
difference, and we do not read the difference as an elasticity in the running
count. The per-request $\tau_2$, which shares no code path with $\tau_1$, gives
$\epsilon = +4.3$ there and rises monotonically across all six points (22.7, 24.5, 25.7, 26.8, 27.6,
29.0\,ms). What the sweep does show is the weaker thing we can support: the flat
regime does not continue indefinitely---\texttt{cap8}+HC realizes 59.7\% of the
bound at the far end of the swept range---but our six points cannot say where
the regime ends.

What the measurement does is rule
one explanation out; it does not establish the opposite
(\S\ref{sec:results}), and the failure is not that this workload denies
admission control its lever, but that we could not build a rule that uses it
well.

What this does and does not establish: the HiCache curve
rests on six configurations and five segments---four computable, the fifth
$\overline{R}$-reversed and therefore given no coefficient---one run each, with
no error bars, stopping at 1.87---so the cliff question is partly answered
rather than merely unasked, since realization does start to fall at the far end
of the swept range. Each run is the first of its configuration, whereas
\S\ref{sec:results}'s frontier uses the
repeats where a pair exists; substituting them moves the first two segments
from $+0.19$ and $+0.46$ to $+0.20$ and $+0.44$, which changes nothing we say,
and it \emph{reverses the third}, because the repeat of \texttt{cap4}+HC sits
at a slightly \emph{higher} running count than \texttt{cap5}+HC, so that
segment's $\overline{R}$ interval runs backwards and its $\epsilon$ comes out
at $-5.2$. We report the third segment and read neither a value nor a location
off it. It is measured on different policies rather than by an
intervention---a cross-policy association, not a transportable causal
derivative. On our data it is also computed from the same wall clock the model
then explains---$\tau_1$ is derived from $W_{\text{busy}}$, so using
it to predict $W_{\text{busy}}$ restates rather than tests---while $\tau_2$,
genuinely independent of the wall clock, measures only decode and cannot be
dropped into a throughput equation that includes prefill.

